\section{POO — Héritage et Polymorphisme}

\subsection{L'Héritage et le Partage de Comportement}

L'héritage permet de créer une hiérarchie de classes en factorisant la logique technique commune dans une classe de base (\textit{Base Class}) et en spécialisant cette logique dans des classes dérivées (\textit{Derived Classes}). En C\#, l'héritage d'implémentation est \textbf{simple} (une classe ne peut hériter que d'une seule classe mère) et s'indique au niveau de la déclaration avec le symbole deux-points (\texttt{:}).

Cette architecture réduit la duplication de code et facilite la maintenance des systèmes à large échelle (comme un système de routage HTTP où tous les routeurs héritent d'une base commune).

\subsection{Méthodes Abstraites, Virtuelles et Scellement}

Le comportement de l'héritage est rigoureusement contrôlé par des modificateurs spécifiques :
\begin{itemize}
    \item \textbf{\texttt{abstract}} : Une classe abstraite est incomplète et ne peut pas être instanciée (l'appel à \texttt{new} est interdit). Elle sert de fondation contractuelle. Une méthode \texttt{abstract} n'a pas de corps et \textbf{doit impérativement} être implémentée par l'enfant.
    \item \textbf{\texttt{virtual}} : Une méthode virtuelle possède une implémentation par défaut exploitable dans la classe mère, mais \textbf{peut} être redéfinie (surchargée) de manière optionnelle par l'enfant via le mot-clé \texttt{override}.
    \item \textbf{\texttt{base}} : Permet à une classe enfant d'appeler explicitement l'implémentation de sa classe mère. Très utilisé dans les constructeurs ou pour enrichir une méthode \texttt{virtual} existante plutôt que de l'écraser complètement.
    \item \textbf{\texttt{sealed}} : Appliqué à une classe, il bloque définitivement l'héritage (aucune autre classe ne peut en dériver). Appliqué à une méthode \texttt{override}, il empêche les enfants sous-jacents de la redéfinir à nouveau. Utilisé pour des raisons de sécurité ou de micro-optimisation de compilation.
\end{itemize}

\subsection{Le Polymorphisme et le Transtypage (\textit{Casting})}

Le \textbf{polymorphisme} est la capacité d'une instance dérivée à être manipulée via une variable typée comme son parent, tout en garantissant que c'est bien le code spécifique à son type réel qui sera exécuté à l'exécution (résolution dynamique). 

Pour vérifier dynamiquement la nature d'un objet en mémoire, le C\# offre des opérateurs de transtypage :
\begin{itemize}
    \item \textbf{\texttt{is}} : Vérifie si un objet est compatible avec un type donné (Pattern Matching).
    \item \textbf{\texttt{as}} : Tente une conversion sécurisée vers un type de référence. Retourne \texttt{null} si la conversion échoue, évitant ainsi un crash applicatif.
\end{itemize}

\subsection{Exemple de Code Commenté}

L'exemple suivant modélise un pipeline de traitement de documents (PDF, XML) démontrant l'héritage et l'exécution polymorphe.

\begin{lstlisting}[caption={Système de traitement documentaire avec polymorphisme}]
// 1. Classe de base abstraite
public abstract class DocumentProcessor
{
    protected string ParserVersion = "1.0";

    // 2. Méthode virtuelle : comportement par défaut surchargeable
    public virtual void Validate(byte[] data)
    {
        if (data == null || data.Length == 0)
            throw new ArgumentNullException(nameof(data), "Le flux binaire est vide.");
    }

    // 3. Méthode abstraite : contrat strict sans corps
    public abstract void Process(byte[] data);
}

// 4. Héritage via le symbole ':'
public class PdfProcessor : DocumentProcessor
{
    // 5. Redéfinition obligatoire de la méthode abstraite
    // 6. 'sealed' bloque la redéfinition pour de futurs enfants
    public sealed override void Process(byte[] data)
    {
        // 7. Appel à l'implémentation parent de validation avec 'base'
        base.Validate(data); 
        Console.WriteLine($"Analyse du flux PDF (Moteur v{ParserVersion}).");
    }
}

public class XmlProcessor : DocumentProcessor
{
    // Redéfinition optionnelle de la méthode virtuelle
    public override void Validate(byte[] data)
    {
        base.Validate(data);
        Console.WriteLine("Vérification approfondie du schéma XSD.");
    }

    public override void Process(byte[] data)
    {
        Console.WriteLine("Désérialisation des noeuds XML.");
    }
}

class Program
{
    static void Main()
    {
        // 8. Polymorphisme : Un tableau parent stockant des enfants hétérogènes
        DocumentProcessor[] processors = {
            new PdfProcessor(),
            new XmlProcessor()
        };

        byte[] fakePayload = new byte[256];

        foreach (var proc in processors)
        {
            // Résolution dynamique de la bonne méthode Process() au runtime
            proc.Process(fakePayload);
            
            // 9. Transtypage avec 'is' (Pattern Matching moderne)
            if (proc is XmlProcessor xmlProc)
            {
                Console.WriteLine("-> Type XML détecté, routage vers l'API externe.");
            }
            
            // Alternative historique avec 'as'
            PdfProcessor pdfProc = proc as PdfProcessor;
            if (pdfProc != null)
            {
                Console.WriteLine("-> Type PDF détecté, génération des vignettes.");
            }
        }
    }
}
\end{lstlisting}

\textbf{Explications ligne par ligne :}
\begin{itemize}
    \item \textbf{Ligne 14} : Définition d'un comportement abstrait. Tout \texttt{DocumentProcessor} doit savoir se traiter (\texttt{Process}), mais l'algorithme exact dépend intimement du format concret.
    \item \textbf{Ligne 22} : \texttt{PdfProcessor} satisfait le contrat en redéfinissant \texttt{Process} avec \texttt{override}. Le mot-clé \texttt{sealed} garantit la fin de la chaîne d'héritage pour ce comportement.
    \item \textbf{Ligne 25} : \texttt{base.Validate()} déclenche la vérification de nullité implémentée en ligne 7, évitant de dupliquer ce code de sécurité métier.
    \item \textbf{Ligne 51} : Le tableau est typé \texttt{DocumentProcessor}, mais il contient en mémoire des instances réelles \texttt{PdfProcessor} et \texttt{XmlProcessor}. C'est la force du couplage lâche.
    \item \textbf{Ligne 60} : L'opérateur \texttt{is} effectue un test de type et affecte directement la variable typée \texttt{xmlProc} si le test réussit, évitant un \textit{cast} manuel redondant.
\end{itemize}

\begin{tcolorbox}[colback=white,colframe=keywordcolor,title=\textbf{Pièges courants \& Erreurs de compilation}]
    \begin{itemize}
        \item \textbf{\texttt{InvalidCastException}} : Tenter de forcer un transtypage de manière brutale avec des parenthèses (ex: \texttt{(PdfProcessor)proc}) lèvera une exception violente si \texttt{proc} est en réalité un \texttt{XmlProcessor}. Il faut \textbf{toujours} privilégier les opérateurs \texttt{is} ou \texttt{as}.
        \item \textbf{Masquage accidentel (\textit{Hiding})} : Si l'enfant déclare une méthode avec la même signature que le parent sans utiliser \texttt{override} (mais avec \texttt{new}), la liaison polymorphique est brisée. Le compilateur déclenchera l'avertissement \texttt{CS0114}.
    \end{itemize}
\end{tcolorbox}

\subsection{Synthèse / À retenir}

\begin{itemize}
    \item \textbf{Factorisation} : Poussez le code commun vers le haut (dans la classe de base) pour respecter le principe DRY (\textit{Don't Repeat Yourself}).
    \item \textbf{Abstract vs Virtual} : Utilisez \texttt{abstract} pour forcer le développeur enfant à écrire du code. Utilisez \texttt{virtual} pour proposer une implémentation par défaut optionnellement modifiable.
    \item \textbf{Couplage Faible} : Typisez vos variables, vos paramètres de méthodes et vos collections avec le type de la classe de base, pas avec le type des enfants, afin de concevoir une architecture ouverte à l'ajout futur de nouveaux comportements.
\end{itemize}
